\documentclass[aspectratio=169,11pt]{beamer}

\usepackage[spanish]{babel}
\usepackage[T1]{fontenc}
\usepackage[utf8]{inputenc}
\usepackage{xcolor}
\usepackage{graphicx}
\usepackage{ragged2e}

\graphicspath{
  {imagnes/}
  {src/}
}

\definecolor{MuseBlue}{HTML}{123C69}
\definecolor{MuseGold}{HTML}{D9A441}
\definecolor{MuseLightBg}{HTML}{FFFFFF}
\definecolor{MuseDarkText}{HTML}{102A43}
\definecolor{MuseSoftBox}{HTML}{F4F7FA}

\usetheme{default}
\setbeamertemplate{navigation symbols}{}
\setbeamertemplate{headline}{}
\setbeamertemplate{footline}{}

\setbeamercolor{background canvas}{bg=MuseLightBg}
\setbeamercolor{normal text}{fg=MuseDarkText,bg=MuseLightBg}
\setbeamercolor{frametitle}{fg=white,bg=MuseBlue}
\setbeamerfont{frametitle}{size=\Large,series=\bfseries}

\newcommand{\uniLogo}{src/logo.png}

\setbeamertemplate{frametitle}{
  \nointerlineskip
  \begin{beamercolorbox}[wd=\paperwidth,ht=1.15cm,dp=0.18cm,leftskip=0.45cm,rightskip=0.35cm]{frametitle}
    \usebeamerfont{frametitle}\insertframetitle
    \hfill
    \includegraphics[height=0.78cm]{\uniLogo}
  \end{beamercolorbox}
}

\setbeamertemplate{title page}{
  \begin{center}
    \vspace{0.8cm}
    \includegraphics[height=2.2cm]{\uniLogo}\par
    \vspace{0.8cm}
    {\LARGE\bfseries\inserttitle\par}
    \vspace{0.35cm}
    {\Large\color{MuseGold}\insertsubtitle\par}
    \vspace{0.9cm}
    {\large\insertauthor\par}
    \vspace{0.2cm}
    {\normalsize\insertinstitute\par}
    \vspace{0.2cm}
    {\normalsize\insertdate\par}
    \vfill
  \end{center}
}

\newcommand{\storyframe}[3]{
  \begin{frame}{#1}
    \begin{columns}[T,totalwidth=\textwidth]
      \begin{column}{0.63\textwidth}
        \begin{center}
          \includegraphics[width=\linewidth,height=0.73\paperheight,keepaspectratio]{#2}
        \end{center}
      \end{column}
      \begin{column}{0.35\textwidth}
        \vspace{0.02\paperheight}
        \colorbox{MuseSoftBox}{%
          \parbox{0.96\linewidth}{%
            \vspace{0.2cm}
            \justifying #3
            \vspace{0.2cm}
          }%
        }
      \end{column}
    \end{columns}
  \end{frame}
}

\title{MuseIQ}
\subtitle{Presentación visual para exposición}
\author{Jhunior Eduardo Guevara Lázaro}
\institute{Universidad Nacional de Ingeniería}
\date{\today}

\begin{document}

\begin{frame}[plain]
  \titlepage
\end{frame}

\storyframe
  {El visitante sí quiere conectar con el museo}
  {Galería vibrante llena de visitantes.png}
  {En 2024, los museos administrados por el Ministerio de Cultura registraron 1\,790\,502 visitas. La demanda presencial sigue siendo alta y confirma que mejorar la experiencia de visita es un problema real y vigente.}

\storyframe
  {Hoy la visita todavía tiene fricciones}
  {sala-de-arte-contemp.png}
  {Una parte importante de la mediación museística sigue dependiendo de cartelas, recorridos lineales y audioguías poco adaptativas. Eso limita personalización, accesibilidad y comprensión situada durante la visita.}

\storyframe
  {Existe una brecha entre interés y mediación}
  {museo futurista con areas contrastantes.png}
  {La brecha central del proyecto está entre una demanda cultural en recuperación y una digitalización todavía parcial de la mediación presencial. MuseIQ se plantea justamente para cerrar esa distancia.}

\storyframe
  {La experiencia puede transformarse}
  {Museo interactivo antes y después.png}
  {La propuesta cambia la lógica de visita: de una experiencia principalmente estática a una experiencia contextual, accesible y más personalizada, apoyada en el propio smartphone del visitante.}

\storyframe
  {MuseIQ convierte el celular en una guía inteligente}
  {Experiencia interactiva en el museo.png}
  {MuseIQ combina smartphone, proximidad BLE, orientación del dispositivo, interacción por voz y recuperación aumentada de información para activar contenido relevante en tiempo real dentro del museo.}

\storyframe
  {La tecnología debe sentirse natural, no invasiva}
  {Museo inteligente con tecnología interactiva.png}
  {El enfoque del proyecto es BYOD: el visitante usa su propio celular. Esto reduce la necesidad de hardware dedicado por institución y facilita una adopción más viable y menos invasiva.}

\storyframe
  {La propuesta integra patrimonio y experiencia digital}
  {Museo con arte y tecnología aumentada.png}
  {La propuesta no se limita a una audioguía ni a un chatbot. Su valor está en integrar contexto espacial, voz, contenido curatorial e inteligencia conversacional en una sola experiencia.}

\storyframe
  {El visitante recibe contenido útil según dónde está}
  {Museo interactivo con arte digitalizado.png}
  {El objetivo técnico no es lograr precisión centimétrica, sino reconocer con confiabilidad suficiente la sala o zona probable del visitante para entregar contenido contextualizado y útil.}

\storyframe
  {La interacción digital ocurre dentro del recorrido}
  {Interacción digital en un museo moderno.png}
  {La experiencia propuesta se integra a la visita presencial. El valor no está en sacar al usuario del museo, sino en acompañarlo digitalmente mientras observa, escucha y pregunta dentro de la sala.}

\storyframe
  {La accesibilidad también es parte del diseño}
  {Conexión cultural a través de la app.png}
  {MuseIQ incorpora TTS y STT como mecanismos de accesibilidad y autonomía. La interacción por voz reduce barreras de lectura sostenida y hace más fluida la experiencia para distintos perfiles de usuario.}

\storyframe
  {No partimos de una idea aislada}
  {Red de museos globales conectados.png}
  {El estado del arte y el mercado internacional muestran una transición clara hacia experiencias móviles, personalizadas y multimodales. MuseIQ se alinea con esa tendencia desde una adaptación al entorno peruano.}

\storyframe
  {El Perú ya tiene condiciones para adoptar esta solución}
  {Espacios culturales conectados con tecnología.png}
  {La tenencia de smartphones en hogares peruanos pasó de 78.0\% en 2019 a 94.8\% en 2024. Ese contexto fortalece la viabilidad de una solución móvil para museos y espacios culturales.}

\storyframe
  {La propuesta también tiene lógica de escalamiento}
  {Visitas modernas en un museo.png}
  {La evaluación económico-financiera de la tesis sugiere condiciones favorables de viabilidad para un despliegue institucional progresivo, apoyado en una arquitectura de bajo costo y servicios recurrentes.}

\storyframe
  {MuseIQ une valor cultural, tecnología y viabilidad}
  {Museo con tecnología y arte.png}
  {El aporte principal de MuseIQ es articular mediación cultural, accesibilidad, contexto espacial e inteligencia conversacional en una propuesta técnicamente defendible y económicamente escalable.}

\begin{frame}{Cierre}
  \begin{center}
    \includegraphics[height=2cm]{\uniLogo}\par
    \vspace{0.5cm}
    {\Large MuseIQ como propuesta de mediación cultural inteligente}

    \vspace{0.35cm}
    {\large Una experiencia más clara, accesible y contextual para el visitante}
  \end{center}
\end{frame}

\end{document}
